\documentclass[11pt]{article}
\usepackage[a4paper,margin=25mm]{geometry}
\usepackage{amsmath,amssymb,amsthm,mathtools}
\usepackage[T1]{fontenc}
\usepackage[colorlinks=true,linkcolor=blue,urlcolor=blue]{hyperref}
\hypersetup{pdftitle={Weighted numerator moments in the quasi-Riemann
argument},pdfauthor={Akash Levy}}
\setlength{\emergencystretch}{2em}
\Urlmuskip=0mu plus 1mu
\newtheorem{theorem}{Theorem}
\newtheorem{lemma}{Lemma}
\newtheorem{proposition}{Proposition}
\newcommand{\Bstar}{\beta_*}
\newcommand{\F}{\mathbb Q(\sqrt{-3})}
\title{Weighted numerator moments in the quasi-Riemann argument}
\author{Akash Levy\thanks{GPT-6 Astra (Ultra) was used to discover the
proof and perform partial formalization; Opus 5.5 Ultracode was used to
complete the formalization and submit the result.}\\[2pt]
{\normalsize GitHub: \texttt{akashlevy}}}
\date{9 October 2026 (PDT)}
\begin{document}
\maketitle

\begin{abstract}
We replace a pointwise numerator bound in the cubic-theta probe by a
weighted fourth-moment estimate for its actual smoothed numerator sum.
Mellin inversion and the Hecke functional equation expose a shorter plain
polynomial, leaving capacity for selected prime factors. A coefficient
norm argument controls the original Euler correction, including its
ramified error terms, without inserting arbitrary row-dependent
coefficients into the moment theorem. Using the original moment parameter
$\kappa=3/4$, the resulting comparison gives
$\beta_*\le10499/12000=0.874916666\ldots$ for the original Hecke family
and hence for every Dirichlet $L$-function: none has a zero with
$\Re s>10499/12000$, the pole of the trivial character at $s=1$ being
excluded. The finite exponent checks are exact rational certificates.
The complete argument is proved in Lean~4 with Mathlib, on top of the
Lean development of the original quasi-Riemann paper. The endpoint
theorems depend only on the axioms \texttt{propext},
\texttt{Classical.choice} and \texttt{Quot.sound}. A Lake build in the
proof package's workspace, with a fresh clone of the pinned fork of the
original development, compiled the 2670 modules of that development
which the proof imports and all 228 modules of the proof. The Lean FRO
Comparator matched the tracker's three template statements for zeta,
for all Dirichlet $L$-functions and for the original Hecke family at
$10499/12000$ against definitions compiled from unmodified upstream
sources; the con-ron and NanoDa kernels and Lean's own kernel accepted
the exported proof. All checks were run by the author on one macOS machine,
the three-kernel check without Comparator's sandbox; they are not a
tracker receipt. No external human peer review has taken place.
\end{abstract}

\section{Statement, scope, and source inputs}

Let $\Bstar$ be the supremum of $1/2$ and the real parts of nontrivial
zeros of the primitive finite-order Hecke $L$-functions over
$\F=\mathbb Q(\zeta_3)$, with arbitrary conductor and height. Principal
poles are excluded.

\begin{theorem}[Improved zero-free half-plane]\label{thm:weighted-main}
\[
 \Bstar\le B:=\frac{10499}{12000}=0.874916666\ldots.
\]
Every finite-order Hecke $L$-function over $\F$, and every Dirichlet
$L$-function including the Riemann zeta function, is therefore zero-free
on the open half-plane $\Re s>B$. Only the pole at $s=1$ of the principal
(trivial) characters is excluded.
\end{theorem}

\paragraph{Lean statements.}
Theorem~\ref{thm:weighted-main} is proved in Lean by the following
unconditional theorems of \nolinkurl{Final.lean}. Lean file names in this
paper are relative to \nolinkurl{formalization/WeightedQRH/} in the
accompanying proof package (Section~\ref{sec:verification}).
\begingroup
\raggedright
\begin{itemize}
\item \nolinkurl{WeightedQRH.beta_le_theta}:
\nolinkurl{HeckeZeroSupremum.beta} $\le10499/12000$.
\item \nolinkurl{WeightedQRH.hecke_nonzero}: for every
\nolinkurl{HeckeFamily.Character} $\chi$ and every $s\in\mathbb C$ with
$10499/12000<\Re s$ and ($s\ne1$ or
$\chi.\mathrm{residue}\ne1$), \nolinkurl{HeckeFamily.LFunction}
$\chi\;s\ne0$.
\item \nolinkurl{WeightedQRH.dirichlet_nonzero}: for every modulus $q$
with \texttt{[NeZero q]}, every complex Dirichlet character $\chi$
modulo $q$ and every $s$ with $10499/12000<\Re s$ and
$\neg(\chi=1\wedge s=1)$, Mathlib's $L$-function
\nolinkurl{DirichletCharacter.LFunction} $\chi\;s$ is nonzero.
\item \nolinkurl{WeightedQRH.zeta_nonzero}: for every $s$ with
$10499/12000<\Re s$ and $s\ne1$, Mathlib's \nolinkurl{riemannZeta} is
analytic at $s$ and \nolinkurl{riemannZeta} $s\ne0$.
\end{itemize}
\endgroup
Here \nolinkurl{HeckeZeroSupremum.beta} is the definition of the original
development \cite{openai,fork}. It is the supremum of $1/2$ and the real
parts of all zeros with $\Re s>0$, other than the pole, of
\nolinkurl{HeckeFamily.LFunction} $\chi$. The supremum runs over every
\nolinkurl{HeckeFamily.Character} $\chi$, given by a nonzero ideal modulus
of the integers of $\F$ and a residue character trivial on the global
units; $\chi$ need not be primitive. This supremum equals $\Bstar$.
Passing to the primitive character removes finitely many Euler factors
$1-\chi(\mathfrak p)N\mathfrak p^{-s}$, whose zeros lie on $\Re s=0$.
The zeros of a primitive $L$-function with $\Re s>0$ are its nontrivial
zeros. The Hecke statements concern exactly these finite-order Hecke
characters of $\F$ in the original definition. They are not a statement
of the generalized Riemann hypothesis, and nothing is claimed for other
fields.

We use the original physical probe, its zero extensions, common masks,
and fixed excluded prime set. No modified zero-free theorem from another
refinement is an input. In particular, the moment parameter remains
$\kappa=3/4$ throughout; no lower-$\kappa$ feedback is used.
The starting theorem $\Bstar\le7/8$ makes the original hypothesis
$\Bstar\le(1+\kappa)/2$ available.

The indispensable source results are:
\begin{itemize}
\item the original zero-free theorem and \texttt{lem:plain}
(fourth moment with short prime factors);
\item \texttt{lem:marked}, \texttt{lem:inverse-amplification}, and
\texttt{prop:detector-witness};
\item \texttt{lem:reflected-energy} and \texttt{prop:probe-gram};
\item \texttt{lem:local-euler} and the exact local formulas underlying
\texttt{prop:probe-errors};
\item the natural-character reflection and \texttt{lem:smooth-calculus};
\item the original contour/tail lemmas and continuation criterion\\
\texttt{lem:continuation-criterion};
\item the transfer theorem \texttt{prop:hecke-dirichlet-transfer}.
\end{itemize}
Labels refer to the \LaTeX{} source of \cite{openai} at the pinned
commit given in the bibliography.
Our new estimate is proved below; it is not assumed as an additional
moment or probe hypothesis. The written proof below cites the source
theorems listed above and accepts them as proved in \cite{openai}. The
Lean proof does not rest on accepting the paper. It uses the Lean
formalization of the corresponding source results in the original
development \cite{fork}, together with Mathlib, and its endpoints depend
only on the standard axioms (Sections~\ref{sec:lean-correspondence}
and~\ref{sec:verification}).

\section{Geometry and the direct bound}

Choose
\begin{equation}\label{eq:new-geometry}
 \ell=\frac{167}{1000},\quad b=\frac{43}{500},\quad
 l_x=\frac{747}{2000},\quad l_y=\frac{919}{2000},\quad
 h=\frac{1587}{2000},\quad \kappa=\frac34.
\end{equation}
Thus $M=l_x+l_y=1-\ell$, $b=l_y-l_x$, and $h=1-l_x+\ell$.
The physical lengths are $X=Z^{l_x}$, $Y=Z^{l_y}$ and
$\prod_iP_i=Z^\ell$, with a fixed even number of disjoint prime slots.
Let $I_{\eta,\mathrm{modified}}(Z)$ be the finite compensated probe
in the source equation \texttt{eq:compensated-probe-definition},
evaluated at these lengths, with the same marked-minus-rescaled
operations and the same original slot windows. Set
\[
 S_i(Z)=\sum_{p\in\mathcal P_i(Z)}W_i(q_p/P_i)q_p^{-5/6},\qquad
 A_\eta(Z)=(-1)^K Z^{-\ell/6}\prod_{i=1}^K S_i(Z),
\]
and define
\[
 J_\eta(Z)=\frac{I_{\eta,\mathrm{modified}}(Z)}{c_S A_\eta(Z)}.
\]
Here $c_S>0$ is the fixed scalar from the source equation
\texttt{eq:shared-principal-data}. The source ray-prime asymptotic gives
$|A_\eta(Z)|\asymp(\log Z)^{-K}$ for sufficiently large $Z$.
Both low and high estimates below concern this same normalized probe.
The normalized principal signal has exponent
\[
 C(s)=s-\frac{4+b}{6}.
\]

For completeness, the direct estimate at these changed lengths follows
from the original estimates without a new cancellation input. For a
rescaled subset of total length $v\in[0,\ell]$, put
$M'=M-2v$ and $\ell'=\ell-v$. Then
$M'+\ell'-1=-3v$. The source reflected-energy proof consequently gives
\[
 E_{\rm ref}\le M'+\frac{(1+3\ell'-2M')_+}{4}+o(1)
 =M'+\frac{(5\ell-1+v)_+}{4}+o(1).
\]
After the Gram bound, Cauchy--Schwarz, tuple counting and rescaling, the
additional exponent is
\[
 -v+\frac{(5\ell-1+v)_+}{8}\le0.
\]
Indeed it is zero at $v=0$ and has slopes $-1$ and $-7/8$, since
$\ell<1/5$. The required Gram domination is
$l_y-\ell-11b/6>0$, and every rescaled length is positive. Therefore
the normalized physical probe satisfies
\begin{equation}\label{eq:new-low}
 |J_\eta(Z)|\ll_{\eta,\epsilon}Z^{l_x/2+b/12+\epsilon}
 =Z^{C(B)+\epsilon},\qquad B=\frac{11}{12}-\frac\ell4.
\end{equation}
All strict scale inequalities used here are checked by the certificate
in Section~\ref{sec:new-certificate}.

\section{Available row counts}\label{sec:row-counts}

Write $U=Z^d$ for a physical row dyad. A dynamic zero bin has
$a=(1+\delta)/2$, with $1/50<\delta\le3/4$ after using $\Bstar\le7/8$.
The floor bin $\delta=1/50$ is handled separately. For the physical main
prime factors, let $q=\delta x$ be their length-weighted mean amplitude,
where $0\le x\le1/2$; assign amplitude zero to an error slot.

Set $\alpha=5/6$, $c=2/9$, and
\begin{align*}
 D_x&=3-(1+4c)x,& P_x&=(2-4cx)(1-x),\\
 J&=(\alpha-\delta)D_x+\delta P_x,&
 R&=1-\delta+\frac{(\alpha-\delta)\delta P_x}{2J}.
\end{align*}
The original selected-witness argument gives
$\#\mathcal C\ll U^{R+\epsilon}(1+T)^A$ for every amplitude/witness
subdivision $\mathcal C$. Here the plain capacity is exactly
$(1-2m)/(6\kappa)=2(1-2m)/9$. At fixed $\kappa=3/4$ there is no
capacity-replacement loss. To see the count directly, the two short
exponents are
\[
 1-\delta\{x+(1-x)r\},\qquad
 1-\delta\{2cx+(2-4cx)(t-r)\}.
\]
Balance their crossing against the long-witness exponent
$1-\delta+(\alpha-\delta)(t-1)$ by choosing
$t=1+\delta P_x/(2J)$. This gives the displayed $R$.
The selected inverse capacity is at most $7/37$, and the plain one at
most $2/27$. Since $\ell/h>7/37$, the original whole-slot selection
and its small capacity decrements apply. The original zero-capacity
arguments are retained. In particular $R\ge1-\delta$.

The row count is first applied to the original amplitude class containing
\emph{all} physical slots. Its full mean $q_{\rm full}$ is at least the
main-slot mean $q$ obtained by setting error-slot amplitudes to zero;
both means are divided by the full slot length. For fixed $\delta$,
$R(\delta,x)$ is nonincreasing on $0\le x\le1/2$. Hence the original
bound with $R(\delta,q_{\rm full}/\delta)$ implies the bound with
$R(\delta,q/\delta)$ used below. This is the Lean theorem
\nolinkurl{WeightedQRH.masked_amplitude_class_count} in
\nolinkurl{HighRows/MaskedCount.lean}. It requires no new moment
assertion for a masked slot family.

\section{The weighted physical numerator estimate}

The useful distinction is between a pointwise bound for
$L(w,\chi_\bullet(u))$ and its smoothed inverse Mellin transform at
the actual scale $Y$. In the latter, reflection produces a plain
polynomial of length at most $U/Y$, rather than the square-root
conductor length used in a generic approximate functional equation.
The unused length supports selected prime factors in the existing fourth
moment. The next proof includes the Euler correction and every error
subset; using only its all-main term would be insufficient.

% Standalone fragment.  Requires amsmath, amssymb, amsthm and hyperref.
% Proof relative to the explicit analytic lemmas of the original manuscript.
\subsection{A weighted moment for the reconstructed physical numerator}
\label{sec:new-weighted-numerator}

The following is a new analytic amendment to the high estimate.  Its
inputs are the exact local identities, smooth reflection and scale
supremum lemmas, and the fourth moment with short prime factors
\textup{(Lemma \texttt{lem:plain})} of the original manuscript.  Those
inputs are imported, rather than reproved here. The coefficient bounds,
literal Mellin reconstruction, original-numerator fourth moment, and
their composition over the amplitude bins are checked in Lean
(Section~\ref{sec:lean-correspondence}); the composition is
\nolinkurl{WeightedQRH.Numerator.source_integrated_amplitude_classes},
on which \nolinkurl{WeightedQRH.Numerator.actual_raw_fullW_input} is
built. The argument below requires coefficient estimates for the Euler
correction, not merely pointwise bounds for an Euler product.

Write $z_0=17/50$, $\kappa=3/4$, let $\theta$ denote the target exponent
($\theta=B$ in Section~\ref{sec:new-certificate}), and assume the
already proved global bound $\beta_*\le7/8$.  Thus the zero-free premise of the imported plain
moment holds.  This choice of $\kappa$ is fixed throughout and uses no
feedback improvement.

\paragraph{The estimate.}
Let $U=Z^d$, $Y=Z^{l_y}$, and let $\mathcal B$ be a fixed dynamic bin of
nonprincipal physical rows $u$ with $q_u\asymp U$.  Put
$\delta=2a-1$, where $51/100\le a\le7/8$. Let $\ell$ be the total physical
prime-slot length.  Suppose
\begin{equation}\label{eq:new-capacity-supply}
 \frac d2<l_y<d,\qquad
 0<\frac{2l_y-d}{6\kappa}\le\ell.
\end{equation}
For each pointwise amplitude and witness subdivision $\mathcal C$ of
this bin, suppose the original row-count argument gives
\[
 \#\mathcal C\ll U^{R+\epsilon}(1+T_1)^A,
 \qquad 0\le q\le\delta/2,
\]
where $q\ell$ is the sum of slot length times amplitude, assigning
amplitude zero to every error slot.  Up to arbitrarily small exponent
losses and a fixed polynomial in the retained heights, its contribution
to the high comparison is bounded by the supremum of
\begin{equation}\label{eq:new-weighted-high-exponent}
 \begin{split}
 E_W(d)={}&a-\theta+h(z_0-1/6)-\frac{l_y}{2}-\frac\ell2+q\ell\\
 &+d\left(\frac{3R}{4}+\frac14-z_0\right)
       -\frac{q(2l_y-d)}{12\kappa}.
 \end{split}
\end{equation}
The supremum is taken after continuing the whole fixed dynamic bin.
Individual amplitude sets are not continued separately.  In particular,
choosing between this estimate and the original high estimate for a
whole bin gives
$\min\{\sup_q E_{\rm old},\sup_q E_W\}$, without requiring an
interchange of the minimum and supremum.

\paragraph{Coefficient norms and the local correction.}
Fix the retained $s,z$ with $\Re s=a+16e$ and $\Re z=z_0$.
Set
\[
 \rho=\delta/4,\qquad \sigma_-=1/2-\rho.
\]
Then $\sigma_->0$ and
$\Re s+\sigma_--1\ge\delta/4>0$.  Both $\Re w=1/2$ and
$\Re w=\sigma_-$ lie in the first Euler region.  For a polynomial in
$Q^{-w}$ define
\[
 \left\|\sum_j b_jQ^{-jw}\right\|_\sigma
       =\sum_j|b_j|Q^{-j\sigma}.
\]
The local denominators $1-D$, $1-R$, and $1-V$ in the original
manuscript are independent of $w$ and uniformly separated from zero.
The factor $H_p$ is a polynomial of degree at most two in $Q^{-w}$ at
$p\nmid u$, and at most one at $p\mid u$.  The exact defect identity
\[
 H_p-1=
 \frac{D(V+W-VW)-VW+(1-V)(1-W)\mathcal E_p}{1-D},
 \qquad \mathcal E_p=P_p^*+D,
\]
therefore gives, by taking absolute coefficients \emph{after} this
cancellation,
\[
 \|H_p-1\|_{\sigma_-}\ll
 \begin{cases}
 Q^{-1-\epsilon_H},&p\nmid u,\\
 Q^{-\epsilon_H},&p\mid u,
 \end{cases}
\]
for a fixed $\epsilon_H>0$.  The same holds at $\sigma=1/2$.
The convergent prime product and the divisor-product estimate imply
that every unselected product, excluding any tuple of distinct slot
labels, has absolute coefficient norm $\ll U^\epsilon$ on both lines.
These assertions are uniform in the retained imaginary parts.

Write
$E_p=G_p+\overline{\chi_p}(u)H_p$.  For an error-slot subset $I$,
the quotient-free part of the correction is
\begin{equation}\label{eq:new-error-coefficient-product}
 F_I(w)=\sum_{(p_i)_{i\in I}}
 \prod_{i\in I}\left[
 W_i(q_{p_i}/P_i)q_{p_i}^{z-1}E_{p_i}\right]
 \prod_{\substack{p\notin S\\p\notin\{p_i:i\in I\}}}H_p.
\end{equation}
All remaining slots retain their actual main factors
$-P_i^{z-1/2}Q_i(u;z)$.  Equation
\eqref{eq:new-error-coefficient-product} follows by expanding the
holomorphic selected-tuple identity.  It does not divide by $H_p$ and
does not alter a main-slot mask or coefficient.

At $p\nmid u$, the exact local error cancellation
\textup{(original equation \texttt{old-eq:5.13b})} yields
\[
 \|E_p\|_\sigma\ll
 Q^{-s_r}+Q^{-6z_0}+Q^{4-5s_r-6z_0}+Q^{1-\sigma-6z_0}
 \qquad(\sigma\in\{1/2,\sigma_-\}).
\]
Each exponent is at most $-1/2$ with a fixed margin.  The ideal count
therefore bounds the sum over a slot, including $Q^{z-1}$, by
$O(P_i^{z_0-1/2})$.

At $p\mid u$, with $j=v_p(u)$, the main factor is zero and
\[
 E_p=G_p=\overline{\eta(p)}Q^s(1-V)\mathcal E_p-Q^{-w}H_p.
\]
After dividing the slot factor by $Q^{z_0-1/2}$, the non-tail boundary
terms of the six-valuation table have the following exponents on
$\Re w=\sigma_-$, at $e=0$:
\[
\begin{array}[b]{c|ccccc}
 j&1&2&3&4&5\\ \hline
 \text{exponents}&-1+\rho&-\delta&-\delta+\rho\ \text{and}\ {-3\delta/2}
   &-3\delta/2+\rho&-5\delta/2
\end{array}.
\]
All are negative; positive $e$ improves them.  The $R,V$ tails and
$-Q^{-w}H_p$ are smaller.  The only exception is the strict term
\[
 -\eta(p)(Q-1)Q^{-s-w},\qquad j\ge2,
\]
whose normalized slot factor costs at most $Q^\rho$ on the left line
and at most one on the central line.  Ramified labels are divisor-many.
For any fixed triangle-expanded tuple, let $J$ consist precisely of
its strict ramified labels.  They are distinct.  Choose the original
small-presentation character, whose ideal coefficients agree with the
physical row. Let $C_u$ be the product of its primitive conductor norm
and full redundant Euler-ideal norm. This combined scale satisfies
the original conductor allocation
\[
 C_u\le C_S\operatorname{rad}(u)
 \le C_*U\prod_{p\in J}Q^{-(j-1)}.
\]
Consequently
$(C_u/(C_*U))^\rho\prod_{p\in J}Q^\rho\le1$.
The single conductor deficit is applied to the whole tuple.

Expand $F_I(w)=\sum_k A_{I,u}(k)q_k^{-w}$ and put
$B_I=\prod_{i\in I}P_i^{z_0-1/2}$.  The preceding estimates prove
\begin{align}
 \sum_k|A_{I,u}(k)|q_k^{-1/2}&\ll U^\epsilon B_I,
 \label{eq:new-central-coefficient-norm}\\
 \left(\frac{C_u}{C_*U}\right)^\rho
 \sum_k|A_{I,u}(k)|q_k^{-1/2+\rho}&\ll U^\epsilon B_I.
 \label{eq:new-left-coefficient-norm}
\end{align}
Combining them gives the useful weighted form
\begin{equation}\label{eq:new-weighted-coefficient-norm}
 \sum_k|A_{I,u}(k)|q_k^{-1/2}
 \max\left(1,\frac{C_uq_k}{C_*U}\right)^\rho
 \ll U^\epsilon B_I.
\end{equation}
This also explains why a strict ramified error does not consume prime
capacity: its forced dilation is paid for by a distinct factor in the
conductor deficit.

\paragraph{Reconstructing and reflecting the numerator.}
For each fixed nonprincipal physical row, move the complete $w$
integral to $\Re w=1/2$.  The numerator is entire; the correction is
holomorphic throughout this move, and the Mellin profile has arbitrary
decay.  The bin itself is independent of $w$.  The absolute coefficient
norm just proved and polynomial strip growth justify termwise Mellin
inversion:
\begin{equation}\label{eq:new-exact-w-reconstruction}
 \begin{split}
 &\frac1{2\pi i}\int_{(1/2)}
 Y^{w-1}\widehat W_1(w)L^S(w,\chi_u)F_I(w)\,dw\\
 &\hspace{8mm}=
 Y^{-1/2}\sum_k A_{I,u}(k)q_k^{-1/2}
              S_{\chi_u}(Y/q_k;W_1).
 \end{split}
\end{equation}
Here $S$ denotes the original centrally normalized plain sum.  Since
$W_1$ has compact support, say in $[r_-,r_+]$, terms with $q_k>r_+Y$
vanish identically before reflection.  Thus all relevant scales lie
in a fixed polynomial range.

The imported smooth functional equation replaces $S_{\chi_u}(X;W_1)$ by reflected sums
with conjugate character and fixed transformed profiles. Restoring
natural zeros follows the original reflection proof. The redundant
Euler restoration need not be supported on fixed $S$: every restoration
subset has norm at most the full redundant ideal norm, and every
smooth-ideal denominator decreases the reflected scale. Consequently
all reflected scales are at most $C_u/X$, with the combined $C_u$
defined above. The original restoration coefficient mass is subpower
in the small modulus. Fixed constants enter $C_*$. No extra weighted
restoration mass or new moving mask is needed.

Use \eqref{eq:new-weighted-coefficient-norm} to remove all row-dependent
coefficients before applying a moment.  It leaves a rowwise supremum
of reflected sums of scale at most
$C_*(U/Y)T$, multiplied by $T^{-\rho}$, where
\[
 T=\max(1,C_uq_k/(C_*U)).
\]
The imported smooth scale-supremum lemma handles the row-dependent
scale using logarithmically many boxes and finitely many derivatives.
Annular decomposition of the reflected profile gives the usual
arbitrarily small length loss, summable lower annuli and arbitrarily
decaying upper tails.  Since $q_k\le r_+Y$, all main reflected scales
are $O(U)$ and the declared lengths stay bounded.  Thus no uniform
moment assertion for unbounded lengths is being used.

\begin{proposition}[Checked original-numerator moment]\label{prop:checked-original-fourth}
Assume $51/100\le\beta_*$. \textup{(}The bound $\beta_*\le7/8$ of
\cite{openai} is available throughout but is not needed here; the Lean
statement assumes only the lower bound.\textup{)}
Fix the original weighted source data, its moment loss $t>0$, a
Schwartz profile $G$, and positive constants $\lambda,\xi,b$.
For sufficiently large $U$, let $\mathcal R$ be any finite set of
nonprincipal physical rows with $U^{1/100}\le q_u\le U$.
For every row choose $0<N_u\le bY$, where $Y>0$, and its actual
combined conductor $C_u$. Suppose, for a common $m\ge0$,
\[
 m+\xi\le2,\qquad
 C_u\le C_S U^m\frac{Y}{N_u}.
\]
Let $T$ be any selected family of original prime slots with nonnegative
widths $r_i$ below the certified mesh
$\operatorname{fineMesh}(2,0,1,3/4,t/4)$, and let
$\Re z_i=17/50$, $|\Im z_i|\le H$, $H\ge0$. Assume
\[
 2(m+\xi)+\frac92\sum_{i\in T}r_i\le1
 \quad\hbox{or}\quad T=\varnothing,
 \qquad
 \left|\prod_{i\in T}Q_i(u;z_i)\right|\le U.
\]
Then there are a fixed integer $J\ge0$ and $C>0$ such that
\[
 \sum_{u\in\mathcal R}
 \left|S_{\chi_u}(Y/N_u;G)\right|^4
 \left|\prod_{i\in T}Q_i(u;z_i)\right|^2
 \le C(1+H)^J U^{1+t+\lambda}.
\]
The constants and threshold precede all choices of rows, scales,
selected slots, and heights. They may depend on the fixed source data,
profile and displayed fixed positive constants.
\end{proposition}
\begin{proof}
Use the actual smooth reflection identity with unit-modulus phase and
cutoff $U^{\xi/4}$. The certified retained-profile theorem bounds the
weighted fourth moment of the retained part by
$C_r(1+H)^J U^{1+t+\lambda}$. It uses the positive-capacity source
moment when $T\ne\varnothing$ and the unrestricted bounded-length
zero-slot source moment otherwise. Both applications preserve the
original physical coefficients and selected factors.

For the discarded part, $C_u\le C_S q_u\le C_SU$ and $N_u\le bY$
give a uniform reflected-scale bound $C_uN_u/Y\le bC_SU$.
The original arbitrary-order discarded-tail estimate therefore gives
any power saving, uniformly over these rowwise choices. The elementary
row count is at most $128U$, and the assumed selected product is at
most $U$. Choosing sufficient tail decay bounds its weighted fourth
moment by a constant $C_d$. Finally,
$|r+d|^4\le8(|r|^4+|d|^4)$, the reflection phase has modulus one,
and $(1+H)^J U^{1+t+\lambda}\ge1$. The result follows with
$C=8(C_r+C_d)$. The Lean theorem allows the more general fixed bound
$C_{\rm slot}U^{q_{\rm tail}}$ on the selected product; the displayed
statement specializes both constants to one.
\end{proof}

\paragraph{Selecting prime factors inside the fourth moment.}
Fix an amplitude set $\mathcal C$ after the full $w$ integration, at
one retained $(s,z)$ and witness tuple.  Put
\[
 L=\ell/d,\qquad n_0=1-l_y/d,
 \qquad z_P=\frac{1-2n_0}{6\kappa}.
\]
For a scale dyad $T=U^v$, allow length $n(v)=n_0+v$ and define
\[
 z(v)=\max\left(0,z_P-\frac{2v}{6\kappa}\right).
\]
The supply condition gives $z_P\le L$.  Select the largest positive
amplitude main slots up to this capacity.  If their total length is
smaller, select all of them.  The original selection lemma, with a
sufficiently fine mesh, gives
\[
 |Q_{\rm selected}|\ge U^{qz(v)-o(1)},
 \qquad |Q_{\rm all}|\le U^{qL+o(1)}
\]
on $\mathcal C$.  Actual selected length never exceeds $z(v)$; no
zero-amplitude factor is required for this lower bound.  The reflected
numerator has character $\overline{\chi_\bullet}(u)$, exactly matching
the physical main slots $\overline{\chi_p}(u)$.  Conjugating the whole
plain/selected-slot product inside its absolute value gives the positive
row orientation in the statement of the imported plain moment; it only
conjugates the profiles and fixed finite-ray coefficients, and preserves
every zero extension.  For sufficiently large $U$ in the moderate
dyads, a physical row cannot induce a character in the fixed exceptional
group $\Theta$, since a good prime of that row ramifies.  The finitely
many unit or fixed-prime-supported exceptions belong to the separate
small-row estimate.

H\"older's inequality and the imported plain moment give
\begin{align*}
 \sum_{u\in\mathcal C}|S_{\bar\chi_u}(n(v))Q_{\rm all}(u)|
 &\ll U^{qL-qz(v)/2+o(1)}(\#\mathcal C)^{3/4}
 \left(\sum_u |S_{\bar\chi_u}(n(v))|^4
                     |Q_{\rm selected}(u)|^2\right)^{1/4}\\
 &\ll U^{qL-qz(v)/2+3R/4+1/4+o(1)}.
\end{align*}
If $z(v)>0$, the capacity equation is $2n(v)+6\kappa z(v)=1$.
If $z(v)=0$, use the original lemma's stronger zero-slot assertion:
it gives $U^{1+o(1)}$ for bounded nonnegative plain lengths, without
a length restriction. This is the certified \texttt{ZeroAt} input,
not the weaker downstream detector estimate with exponent
$\max(1,2n)$. Small slack is reserved for annuli and scale boxes.
The derivative versions used by the scale supremum have the same
exponent and a fixed height cost.

The coefficient weight adds $-\rho v$.  Before the capacity is
exhausted, the exponent relative to $v=0$ is
\[
 \left(\frac q{6\kappa}-\rho\right)v
 \le-\frac{5\delta}{36}v,
\]
because $\kappa=3/4$, $q\le\delta/2$, and $\rho=\delta/4$.
After capacity is exhausted, its slope is $-\rho$. Thus every
dilation dyad is controlled by $v=0$, with room for the small fixed
losses. All error subsets satisfy the same estimate; no lower bound
for an error factor is used.

Restoring the main-slot normalizers, the $w$-integrated bin is bounded
by
\[
 Y^{-1/2}Z^{\ell(z_0-1/2)+q\ell}
 U^{3R/4+1/4-qz_P/2+o(1)}.
\]
The remaining $q_u^{-z}$ factor, reciprocal-denominator estimate and
$X,Z$ Mellin weights are those of the original high-bin proof.
They give exactly \eqref{eq:new-weighted-high-exponent}.  The finite
amplitude, witness and error subdivisions are used only in this
estimate of the already continued whole bin.  Rapid external Mellin
decay and the fixed height orders treat the joins and discarded
coordinates with the original cumulative height allocation.

\paragraph{Dependence on the row scale.}
In the positive-capacity range the coefficient of $d$ in
\eqref{eq:new-weighted-high-exponent} is
\[
 \frac{3R}{4}+\frac14-z_0+\frac q{12\kappa}.
\]
For the original selected count $R\ge1-\delta$ and
$\delta\le5/6$, this is at least $7/200>0$.  Thus $d=h$ controls
the moderate row range whenever \eqref{eq:new-capacity-supply} holds
throughout that range.  The small and far row estimates can remain
those of the original proof.

Finally, at fixed $R,q$ the difference from the original high exponent
is
\[
 E_W-E_{\rm old}
 =\frac{\delta l_y}{2}
   +d\left(\frac{1-R}{4}-\frac\delta2\right)
   -\frac{q(2l_y-d)}{12\kappa}.
\]
The final term is the additional saving supplied by the reconstructed
physical numerator and the existing prime-weighted fourth moment.
It is independent of any parameter-feedback argument.


\section{Exact exponent comparison}\label{sec:new-certificate}

Let $z_c=17/50$. The source pointwise numerator argument gives, relative
to $C(B)$,
\begin{equation}\label{eq:old-high-new-geometry}
 E_{\rm old}(d)=a-B+h(z_c-1/6)-a l_y-\ell/2+q\ell
                  +d(R+\delta/2-z_c).
\end{equation}
For $1/2\le d\le h$, we have $l_y<d<2l_y$. The new weighted estimate
instead gives $E_{\rm wt}=E_W$, that is \eqref{eq:new-weighted-high-exponent}
with $\theta=B$ and $z_0=z_c$:
\begin{equation}\label{eq:weighted-high}
 E_{\rm wt}(d)=a-B+h(z_c-1/6)-l_y/2-\ell/2+q\ell
       +d\left(\frac{3R+1}{4}-z_c\right)
       -\frac{q(2l_y-d)}{12\kappa}.
\end{equation}
The last term is the additional saving from the prime factors in the
numerator moment. At a fixed zero bin we choose the source contour when
$\delta\le1/3$, and the new numerator treatment when $\delta\ge1/3$.
This choice is made for the original dynamic-bin sum before any pointwise
amplitude subdivision. It does not choose a contour according to $x$.

\begin{lemma}\label{lem:new-certificate}
For the geometry \eqref{eq:new-geometry},
\begin{align*}
 E_{\rm old}(h)&\le-10^{-4}
 &&(0\le\delta\le1/3,\ 0\le x\le1/2),\\
 E_{\rm wt}(h)&\le-10^{-4}
 &&(1/3\le\delta\le3/4,\ 0\le x\le1/2).
\end{align*}
All prime-supply, rescaling, remaining row-range, and Euler-region
inequalities needed in this proof hold with the stated strict margins.
\end{lemma}
\begin{proof}
Here is an exact, reproducible finite certificate. Put
$A_P=(\alpha-\delta)\delta P_x$, so that the first identity below holds,
and form the polynomial $T$:
\[
 RJ=(1-\delta)J+\frac{A_P}2,\qquad
 T=\frac{(1+\delta)(1-l_y)}2-\frac\ell2
       +q\ell+\frac{h\delta}{2}-\frac h6-B.
\]
Then the two endpoint numerators are
\begin{align*}
 J E_{\rm old}(h)&=JT+h(RJ),\\
 J E_{\rm wt}(h)&=J E_{\rm old}(h)
   +\frac h4[J-RJ]+\frac{\delta(l_y-h)}2J
   -\frac{q(2l_y-h)}{12\kappa}J.
\end{align*}
All their coefficients are rational and $J>0$ on the full rectangle.
The script
\nolinkurl{evidence/certificate/weighted_numerator_certificate.py} in the
proof package proves that $-J(E+10^{-4})\ge0$, using the convex-hull property
of the Bernstein basis: one rectangle suffices for the first branch,
and nine exact rational subrectangles suffice for the second.
The endpoints and lower Bernstein coefficients are stored in
\nolinkurl{evidence/certificate/weighted_numerator_certificate.json}.
No floating-point sign decision is used. Lean proves the two branch
inequalities separately, in the extended form used by the proof
(Section~\ref{sec:lean-correspondence}).

The same script checks 32 auxiliary rational comparisons, including
\[
 \frac\ell h>\frac7{37},\qquad 5\ell-h>0,\qquad
 2l_y-h>0,\qquad l_y<\frac12,
 \qquad \ell>\frac{2l_y-1/2}{6\kappa}.
\]
The last inequality supplies every numerator capacity on $d\in[1/2,h]$.
For $d\le1/2$, use the original no-slot count
$R=76/75-2\delta/3$; its exponent at $d=1/2$ is strictly negative
at $\delta=1/50$ and $\delta=3/4$, hence throughout by affinity.
The floor-bin estimate with $R=1$, $q=1/100$ is strictly negative
at $d=h$.
For $d\le1/100$ the original absolute small-row comparison has saving
\[
 -\{h(z_c-1/6)-l_y/2+63/5000\}>0
\]
directly relative to $C(\Bstar)$.

Both endpoint estimates control the full main row interval. The source
frequency slope $R+\delta/2-z_c$ is positive. The new slope is
\[
 \frac{3R+1}{4}-z_c+\frac q{12\kappa}
 \ge \frac{7}{16}-\frac{17}{50}>0,
\]
using $\delta\le3/4$ and $R\ge1-\delta$. All these slopes have a
fixed upper bound, so a sufficiently small fixed extension above $h$
preserves the strict endpoint and prime-supply margins.
\end{proof}

\section{Euler regions, parameter choices, and continuation}

The principal comparison uses the same residue as the source. Its second
Euler region extends to $\Re s\ge B$ with $\Re w\ge19/20$ and
$\Re z\ge33/200$: the good-prime defect exponents
$4-6\Re s-6\Re z$ and $1-\Re s-\Re w-6\Re z$ remain strictly below
$-1$. The four principal approximation exponents
\[
 -\Re s,\quad -6\Re z,\quad 4-5\Re s-6\Re z,
 \quad 1-\Re w-6\Re z
\]
are at most $-B$. All local denominators stay uniformly separated from
zero. Thus the principal multiplier is $-1+O(q_p^{-B})$; the normalized
correction $H_\eta(s)=\mathcal H_{\eta,1}(s,1,1/6)$ is holomorphic
on $\Re s>B$ and can be made
within $1/2$ of one by increasing the fixed excluded-prime cutoff.
The principal normalization has a subpower reciprocal.
The source contour lemmas with a printed $7/8$ lower endpoint extend
to $B$ by precisely these local checks. On the absolute small-row line,
choose their first-region buffer below $B-1/2$.

Suppose for contradiction that $\Delta=\Bstar-B>0$.
First choose a fixed row extension $\zeta>0$ smaller than the endpoint,
supply, and $2l_y-h$ margins. Choose the moment, capacity, reflection
truncation, and rounding losses relative to $10^{-4}$ and $\Delta$.
Both the original row counts and the new numerator moment use the
original uniform mesh assertion, so one common mesh suffices.
Next choose the fixed even number of disjoint slots and their windows.
Only after this choice fix the principal approximation loss
$0<\mu<B\min_i\ell_i$. Choose the remaining bin widths, contour buffers,
normalization losses and detector losses small relative to the already
fixed margins and $\mu$. The uniform coefficient norm and reflected
scale arguments in the new lemma consume arbitrarily small losses;
they do not change this order of choices.

All these real choices precede the target character. For each target
choose its fixed arithmetic data, cutoff, and finite internal seminorm
orders. Finally use the source's late height closure: choose the positive
height exponent, then the external tail order, then the threshold.
The new moment/reflection estimates have the same finite-seminorm and
polynomial-height form as the cited source lemmas. To ensure that the
new Mellin identity uses the full $w$ line, the checked transport first
uses the original physical-row-to-central-cube theorem, with
$\Re s=a+16e$, $\Re w=1-a-6e$ and $\Re z=z_c$.
Extend only the truncated $w$ interval to its entire vertical line;
the original polynomial kernel bounds and arbitrary Mellin decay bound
the added tail, uniformly on the retained $s,z$ rectangle.
Then move this complete $w$ line to $1/2$. Along the move
$\Re(s+w)\ge1+10e$, the numerator is entire, and the denominator is
independent of $w$. The first Euler-region bounds and Mellin decay
control the horizontal joins. Both tail errors have arbitrary power
saving, including after summing actual rows and prime tuples.
This does not extend the buffered reciprocal beyond its authorized
$s$ height range. Fubini then gives the full numerator integral inside
the two remaining Mellin integrals. Their Gaussian and radial-Mellin
profile has uniformly bounded moments, so integrating the numerator
estimate introduces no additional contour-height loss. No amplitude
indicator is analytically continued.

Rows above $Z^{h+\zeta}$ use the unchanged absolute large-row contour;
small and intermediate rows use the comparisons certified above. The
normalized probe therefore obeys, for some $0<\omega<\Delta$ and
$\sigma>0$ independent of the target,
\[
 |J_\eta(Z)|\ll_\eta Z^{C(B)+\omega},\qquad
 |J_\eta(Z)-f_\eta(Z)|\ll_\eta Z^{C(\Bstar)-\sigma},
\]
where
\[
 f_\eta(Z)=\frac1{2\pi i}\int_{(2)}
 Z^{C(s)}e^{(s-5/6)^2}\frac{H_\eta(s)}{L_F^S(s,\eta)}\,ds.
\]
These are the source continuation criterion's hypotheses at $B$,
contradicting $\Bstar>B$. This proves the Hecke conclusion.
Restoring finite Euler factors introduces no zeros in this positive
half-plane. The norm-pullback factorization
$L_F(s,\chi\circ N)=L(s,\chi)L(s,\chi\chi_{-3})$ gives the Dirichlet
conclusion exactly as in the source transfer theorem. \qed

\section{Public comparison and limits}

The strongest comparable public unconditional claim located by a dated
search, completed on 10 October 2026 at 00:16 UTC, is the unique root
$B_{\rm public}$ in $(0.8749570194,0.8749570195)$ of
\[
 7884B_{\rm public}^3-18819B_{\rm public}^2
 +14643B_{\rm public}-3686=0.
\]
It is approximately $0.87495701942009894613$. The endpoint appears in
written manuscripts of Shunia with GPT-6 Astra \cite{shunia} and of Guo
\cite{guo}. A Lean package for it by Nielstron is in an open, unmerged
pull request to the QuasiRiemannTracker \cite{nielstron}; its author
reports local Comparator and Lean/NanoDa/con-ron checks of seven targets.
The best value then catalogued by the tracker, with status
framework-verified, is the rational
$874957019420098946128604623/10^{27}$ \cite{tracker}. These reports were
read, not rerun. The present $B=10499/12000$ is smaller than
$B_{\rm public}$ by about $4.0353\cdot10^{-5}$. This is plain arithmetic:
\[
 \frac{10499}{12000}=0.8749166\ldots<0.8749570194<B_{\rm public},
\]
where $0.8749570194$ is the lower endpoint of the interval above. The
comparison with the tracker's rational,
$10499/12000<874957019420098946128604623/10^{27}$, is the Lean theorem
\mbox{\nolinkurl{WeightedQRH.theta_lt_public}} in
\nolinkurl{Parameters.lean}, a module in the endpoints' import closure;
the proof itself does not use it. The only smaller public value
found, Guo's $328/375$ \cite{guo}, is explicitly conditional on an
unproved mean-square hypothesis (Hypothesis~11.1 there). The search is
dated and is not an exhaustive novelty or priority certificate. The
public records were consulted as benchmarks only; no outside refinement
proof code was fetched, used or reproduced here.
The new ingredient is the selected-prime saving in a moment of the
\emph{physical numerator}, which is absent when one bounds that
numerator pointwise and uses moments only to count zero witnesses.
The initial unweighted numerator fourth-moment variant did not improve
the global comparison; retaining the actual reflected length is essential.

The zero-slot argument uses the stronger
certified source theorem, with no restriction on bounded plain lengths;
a weaker detector estimate would not suffice there. Exact exponent
certificates alone do not verify the analytic deduction; that deduction
is checked by the Lean development of
Sections~\ref{sec:lean-correspondence} and~\ref{sec:verification}.
No external human peer review is claimed.

This change invalidates the applicability of the old
row-count/pointwise-numerator optimization barrier to the new high bound. It does not remove
the unchanged direct estimate's obstruction to GRH. Let $g$ be the
extra Gram exponent and $e_{\rm row}$ the extra completed-row energy
exponent, distinct from the small contour buffer used above. If they
satisfy
\[
 g\ge b/6,\quad e_{\rm row}\ge0,\quad
 e_{\rm row}\ge(1+3\ell-2M)/4,\quad e_{\rm row}\ge M+\ell-1,
\]
then the low comparison requires
$B\ge5/6+M/12-\ell/6+e_{\rm row}/2$, and
\begin{align*}
 \frac M{12}-\frac\ell6+\frac{e_{\rm row}}2-\frac1{30}
 ={}&\frac15\left[e_{\rm row}-\frac{1+3\ell-2M}{4}\right]\\
 &+\frac1{60}[e_{\rm row}-(M+\ell-1)]
       +\frac{17e_{\rm row}}{60}\ge0.
\end{align*}
Thus those particular low estimates still force $B\ge13/15$.
Reaching $1/2$ needs a change on that side as well, such as a proved
arithmetic cancellation estimate that survives the reflection phase.
This is a limitation of the stated low estimates only, not a lower bound
for actual zeros or a theorem excluding other methods.

\section{Correspondence with the checked Lean theorems}
\label{sec:lean-correspondence}

\begingroup
\raggedright
The new modules are under \nolinkurl{formalization/WeightedQRH/} in the
proof package. They import the original Lean development \cite{fork} at
the pin given in Section~\ref{sec:verification}. Some of them, for
example the low-side detector and reflection modules, adapt original
modules to the changed lengths and record that origin in their headers.
The following list maps the steps of the written proof to the theorems
composed by the endpoints of \nolinkurl{Final.lean} and
\nolinkurl{Statements.lean}. Apart from the standalone arithmetic barrier
\nolinkurl{physical_low_barrier}, every theorem named in the list lies on
the proof path of one of these endpoints. The names are the actual
theorem declarations in the Lean sources, not additional hypotheses
standing in for the corresponding proofs. Names
from \nolinkurl{HighRows/SourceProbe.lean},
\nolinkurl{HighRows/WeightedSourceRawCells.lean} and
\nolinkurl{HighRows/ConditionalFinal.lean} lie in the namespace
\nolinkurl{OAI.SevenEighths.WeightedHighFinalAssembly}; the others lie in
\nolinkurl{WeightedQRH} or one of its subnamespaces.

\begin{itemize}
\item The low-side comparison at the changed lengths is the actual
\nolinkurl{sourceProbe_low} theorem in \nolinkurl{HighRows/SourceProbe.lean}.
The arithmetic geometry and scoped $13/15$ obstruction are in
\nolinkurl{Parameters.lean}; the obstruction is
\nolinkurl{physical_low_barrier}, which the proof does not use.
\item The polynomial certificates and their real-rectangle implications
are in \nolinkurl{CertificateBernstein.lean} and
\nolinkurl{ExponentCertificate.lean}. In particular,
\nolinkurl{weighted_extended_bound} (with \nolinkurl{old_extended_bound}
for the source branch) includes the fixed row extension and
the $1/20000$ margin. These arithmetic theorems alone imply no
zero-free result.
\item The masked row count of Section~\ref{sec:row-counts} is
\nolinkurl{masked_amplitude_class_count} in
\nolinkurl{HighRows/MaskedCount.lean}.
\item The combined primitive and redundant conductor bound is
\nolinkurl{physical_reflection_data_exists} in
\nolinkurl{NumeratorReflectionData.lean}. The coefficient mass
\eqref{eq:new-weighted-coefficient-norm} is
\nolinkurl{actual_physical_label_mass} in
\nolinkurl{NumeratorPhysicalPacketMass.lean}; it is stated for a general
scale $C_u$ and is applied with this combined conductor.
\item On the line $\Re w=1/2$, each error-label numerator of the
complete $w$ integral is identified with its coefficient packet, as in
\eqref{eq:new-exact-w-reconstruction}, by
\nolinkurl{errorLabelNumerator_reconstruction} in
\nolinkurl{NumeratorLabelReconstruction.lean}. The theorems
\nolinkurl{normalized_prime_tuple_reconstruction} and
\nolinkurl{integratedDyadValue_reconstruction} in
\nolinkurl{NumeratorReconstruction.lean} give the complete finite row and
error-subset decomposition.
\item Proposition~\ref{prop:checked-original-fourth} is
\nolinkurl{certified_original_numerator_fourth} in
\nolinkurl{NumeratorOriginalFourth.lean}. It is proved using the actual
certified source moments, rather than assuming a new reflection or
fourth-moment estimate.
\item The full-$w$ contour transport and uniform remaining Mellin
integration are in \nolinkurl{FullWContour.lean} through
\nolinkurl{FullWBound.lean}, including \nolinkurl{SourceSZBound.lean}; the
transfer from the central cube to the full $w$ line is in
\nolinkurl{HighRows/CubeFullWTransfer.lean} and
\nolinkurl{HighRows/CubeFullWNormalized.lean}. The endpoint integration
theorem is \nolinkurl{FullWContour.finiteFullW_uniform_norm_bound}.
\item Final assembly.
\nolinkurl{Numerator.source_integrated_amplitude_classes} in
\nolinkurl{NumeratorSourceCollect.lean} collects the actual physical error
masks and the original full amplitude classes over all rows of a dyad,
with one constant fixed before the row scale and all target, height and
contour choices. Combined with the masked count, the scalar, conductor
and prime-supply estimates, it proves
\nolinkurl{Numerator.actual_raw_fullW_input}, of type
\nolinkurl{RawFullWAnalyticInput}, in
\nolinkurl{NumeratorActualSource.lean}: the raw full-$w$ bound for every
source row dyad. The theorem \nolinkurl{fullW_input_of_raw_source} in
\nolinkurl{HighRows/WeightedSourceRawCells.lean} adds the source geometry,
height absorption, remaining Mellin integration and normalization, and
yields \nolinkurl{FullWAnalyticInput}. This gives
\nolinkurl{WeightedQRH.fullW_analytic_input} in \nolinkurl{Final.lean}.
\item The implications \nolinkurl{beta_le_theta_of_fullW_input},
\nolinkurl{hecke_of_fullW_input}, \nolinkurl{dirichlet_of_fullW_input} and
\nolinkurl{zeta_of_fullW_input} in
\nolinkurl{HighRows/ConditionalFinal.lean} pass from
\nolinkurl{FullWAnalyticInput} to the common probe bounds
(\nolinkurl{common_probe_of_fullW_input}) and through the continuation
argument of \nolinkurl{Continuation.lean} to the zero-free conclusions.
Applied to \nolinkurl{WeightedQRH.fullW_analytic_input}, they give the
unconditional endpoints \nolinkurl{beta_le_theta},
\nolinkurl{hecke_nonzero}, \nolinkurl{dirichlet_nonzero} and
\nolinkurl{zeta_nonzero} of \nolinkurl{Final.lean}. No hypothesis
remains. The docstring of \nolinkurl{HighRows/ConditionalFinal.lean}
predates \nolinkurl{Final.lean} and describes that module in isolation,
where \nolinkurl{FullWAnalyticInput} is still a hypothesis; the input is
discharged in \nolinkurl{Final.lean}.
\item \nolinkurl{Statements.lean} restates the endpoints exactly in the
form of the Comparator challenges of Section~\ref{sec:verification}:
\nolinkurl{OAI.riemannZeta_ne_zero_of_10499_12000_lt_re},
\nolinkurl{OAI.DirichletCharacter.LFunction_ne_zero_of_10499_12000_lt_re}
and
\nolinkurl{OAI.SevenEighths.HeckeFamily.LFunction_ne_zero_of_10499_12000_lt_re}.
\end{itemize}

The endpoint audit \nolinkurl{FinalAxioms.lean} prints the statements
and the axioms of \nolinkurl{actual_raw_fullW_input},
\nolinkurl{fullW_analytic_input}, the four endpoints of
\nolinkurl{Final.lean} and the three statements of
\nolinkurl{Statements.lean}, and the axioms (not the statement) of
\nolinkurl{source_integrated_amplitude_classes}. Each of these ten
declarations depends only on \nolinkurl{propext},
\nolinkurl{Classical.choice} and \nolinkurl{Quot.sound}. No analytic
axiom is introduced anywhere.
\par
\endgroup

\section{Formal verification and reproducibility}\label{sec:verification}

\begingroup
\raggedright
\paragraph{Proof package.}
The Lean proof accompanies this manuscript in the proof package
\nolinkurl{proofs/akashlevy-20261009-weighted-numerator/} prepared for
the QuasiRiemannTracker \cite{tracker}. Paths in this section are
relative to that directory; Lean file names remain relative to
\nolinkurl{formalization/WeightedQRH/}. The package contains:
\begin{itemize}
\item \nolinkurl{formalization/}: a Lake workspace with
\nolinkurl{lakefile.toml}, \nolinkurl{lean-toolchain},
\nolinkurl{lake-manifest.json}, the scripts \nolinkurl{setup.sh} and
\nolinkurl{run_comparator.sh}, the 228 Lean modules of the proof in
\nolinkurl{WeightedQRH/}, the three Comparator challenges with their
configurations in \nolinkurl{ComparatorChallenges/}, and the
three-kernel driver \nolinkurl{kernels/judge_three_kernels.py}. The 228
modules are the 227 local modules in the import closure of the endpoint
modules \nolinkurl{Final} and \nolinkurl{Statements}, plus the audit
module \nolinkurl{FinalAxioms};
\item \nolinkurl{manuscript/}: the \LaTeX{} sources of this paper;
\item \nolinkurl{evidence/certificate/}: the exact exponent certificate
\nolinkurl{weighted_numerator_certificate.py} and its output
\nolinkurl{weighted_numerator_certificate.json};
\item \nolinkurl{evidence/kernels/}: the report, judge log,
configuration, export targets, control logs, source audit, tool pins,
compile and export logs and the Lean sources \nolinkurl{src/Challenge.lean}
and \nolinkurl{src/Solution.lean} of the three-kernel run, and, in
\nolinkurl{pristine-workspace/}, the Lake configuration and build log of
its challenge workspace;
\item \nolinkurl{evidence/lake/}: the logs of the Lake build, of its
rehash check and of the Comparator runs in that workspace;
\item \nolinkurl{evidence/isolated-rebuild/} and
\nolinkurl{evidence/comparator/}: the record of an earlier isolated
rebuild and its Comparator runs;
\item \nolinkurl{evidence/source-scan.txt}, and the license and notice
files \nolinkurl{LICENSE}, \nolinkurl{NOTICE},
\nolinkurl{UPSTREAM-NOTICE} and \nolinkurl{third-party-notices/}.
\end{itemize}

\paragraph{Toolchain and source pins.}
Lean v4.34.1 (commit
\nolinkurl{5045d0056413266e57c625dcd7c365b10e377c52}); Mathlib commit
\nolinkurl{d13f23b723b8a846827a245b89c10fc7d3f11612} \cite{mathlib};
PrimeNumberTheoremAnd commit
\nolinkurl{c39a751132c88b6e8080b74c74023fd95b3d8be0} and
rellich-kondrachov commit
\nolinkurl{70f85d4c1bf99c6e7d61e8be4daa6f3664d08d23}, each with the
Lean 4.34.1 compatibility patch \nolinkurl{lean/patches/*-lean4341.patch}
of the original development. All modules are compiled with
\nolinkurl{autoImplicit} disabled. The original Lean development is
\nolinkurl{github.com/akashlevy/math} commit
\nolinkurl{c388e88e6519d139ad2ca5571ef0c7fbf0e245b5}, branch
\nolinkurl{qrh-23-24-pruned} \cite{fork}. This fork commit is
\nolinkurl{github.com/openai/math} commit
\nolinkurl{fd4aeeb2ee4fc729c18d98444fed42fd0529eeeb} with the declarations
that no QRH result uses pruned in place, plus 23/24 modules
(\nolinkurl{OAI.NumberTheory.DirichletL.QRH*}) that the present proof
does not import. The \LaTeX{} source of \cite{openai} is byte-identical
at both commits. One adapted module,
\nolinkurl{LowDetector/LowOriginalPools.lean}, previously imported an
original module that the pruning deleted; it now imports that module's
own imports (\nolinkurl{OAI.NumberTheory.DirichletL.Detector.LowSlotScales}
and \nolinkurl{OAI.NumberTheory.DirichletL.Detector.RayPoolGood}) instead.
The three-kernel check additionally uses the official Lean v4.35.0-rc2
release for \nolinkurl{darwin_aarch64}. The Challenge and Solution
exports of that check, and the Comparator runs in the Lake workspace and
on the isolated rebuild, use lean4export commit \nolinkurl{076e8e57}
built with Lean 4.34.1; the preflight controls of the three-kernel check
were compiled and exported with the \nolinkurl{lean} and
\nolinkurl{leanexport} of the v4.35.0-rc2 judge toolchain.

\paragraph{Source scan.}
The file \nolinkurl{evidence/source-scan.txt} records a scan, with
comments stripped, of the 231 Lean files in \nolinkurl{formalization/}
(228 in \nolinkurl{WeightedQRH/}, 3 in \nolinkurl{ComparatorChallenges/})
for \nolinkurl{sorry}, \nolinkurl{admit}, \nolinkurl{axiom},
\nolinkurl{unsafe}, \nolinkurl{implemented_by}, \nolinkurl{extern},
\nolinkurl{skipKernelTC}, \nolinkurl{native_decide},
\nolinkurl{ofReduceBool}, \nolinkurl{trustCompiler},
\nolinkurl{run_cmd}, \nolinkurl{run_elab}, \nolinkurl{run_meta},
\nolinkurl{addDecl}, \nolinkurl{elab}, \nolinkurl{macro_rules},
\nolinkurl{macro}, \nolinkurl{syntax}, \nolinkurl{opaque},
\nolinkurl{partial}, \nolinkurl{csimp}, \nolinkurl{initialize} and
\texttt{\#eval}. No module in \nolinkurl{WeightedQRH/} contains any of
them. The three challenge modules contain \nolinkurl{sorry} by design,
since their statements are the targets. The only options set are
\nolinkurl{maxHeartbeats} and \nolinkurl{maxRecDepth}. The scan does not
cover the original development or Mathlib; the axiom checks and kernel
replays below cover every declaration that the endpoints use.

\paragraph{Lake build from public sources.}
The script \nolinkurl{formalization/setup.sh} clones
\nolinkurl{github.com/akashlevy/math} at \nolinkurl{c388e88e6} into
\nolinkurl{formalization/math/} (a sparse checkout of
\nolinkurl{lean/OAI/NumberTheory/DirichletL/}, \nolinkurl{lean/patches/}
and \nolinkurl{LICENSE}), lets Lake fetch the dependencies pinned in
\nolinkurl{lake-manifest.json}, applies the two compatibility patches,
downloads Mathlib's build cache, and runs
\texttt{lake build WeightedQRH ComparatorChallenges}. In the recorded
run, \nolinkurl{formalization/math/} was a fresh sparse clone of the
fork at that commit. The dependency directory \nolinkurl{.lake/packages}
was instead an existing checkout of the pinned dependencies with the two
patches applied, and the build directory had been seeded with a copy of
an earlier build of the original development. The build reported
``Build completed successfully (11847 jobs)'' after 39 minutes
(\nolinkurl{evidence/lake/lake-build.log}). Its log lists as built in
that run the 2670 modules of the original development in the import
closure of the proof, all 228 proof modules and the 3 challenge modules; Mathlib and the other dependencies entered as
prebuilt objects of the existing checkout. A subsequent
\texttt{lake build -{}-no-build -{}-rehash WeightedQRH ComparatorChallenges},
which rehashes the sources and builds nothing, reported ``All targets
up-to-date (11847 jobs)'' (\nolinkurl{evidence/lake/lake-no-build-rehash.log}).
The output of \nolinkurl{FinalAxioms.lean} in the build log shows only
\nolinkurl{propext}, \nolinkurl{Classical.choice} and
\nolinkurl{Quot.sound} for each of the ten audited declarations
(Section~\ref{sec:lean-correspondence}). Every theorem on the
endpoints' proof paths therefore uses at most these axioms. The recorded
run did not exercise the dependency fetch, patch and cache steps of
\nolinkurl{setup.sh}; its patch step was tested separately on a fresh
clone of rellich-kondrachov at the pinned revision.

\paragraph{Comparator in the Lake workspace.}
The Lean FRO Comparator \cite{comparator} exports a solution theorem with
all its dependencies. It checks that the statement and the definitions it
uses match the challenge, checks the axioms against the permitted list
\nolinkurl{propext}, \nolinkurl{Quot.sound}, \nolinkurl{Classical.choice},
and replays the whole export in the kernel. The challenges in
\nolinkurl{formalization/ComparatorChallenges/} import only Mathlib.
\nolinkurl{WeightedQRHZeta} and \nolinkurl{WeightedQRHDirichlet} use only
Mathlib definitions. \nolinkurl{WeightedQRHHecke} is OpenAI's upstream
\nolinkurl{HeckeSevenEighths} challenge with only the theorem name and the
constant changed. Their solution module is
\nolinkurl{WeightedQRH.Statements}. The script
\nolinkurl{formalization/run_comparator.sh} runs Comparator (commit
\nolinkurl{d03acab1}, default Lean kernel, NanoDa disabled) under
\texttt{lake env} in the workspace above; Comparator's own
\texttt{lake build} step ran as usual. For each of the three challenges,
the log \nolinkurl{evidence/lake/comparator-lake-WeightedQRH*.log} shows
``Lean default kernel accepts the solution'' and ``Your solution is
okay!'', with exit status 0. Like OpenAI's upstream $7/8$ challenge, the
zeta challenge has no hypothesis $s\ne1$. It holds at $s=1$ only because
Mathlib's total \nolinkurl{riemannZeta} assigns a nonzero value at the
pole, which is not an analytic statement there. The analytic statement is
\nolinkurl{WeightedQRH.zeta_nonzero}.

\paragraph{Three-kernel check of the tracker template.}
The driver \nolinkurl{formalization/kernels/judge_three_kernels.py} was
written for this run. It follows the protocol of the tracker's earlier
kernel dossiers (PalomarSubmission \nolinkurl{d4e41c1}, driving
\texttt{lake comparator} from Lean v4.35.0-rc2) and re-implements their
preflight controls, ill-typed export and export targets; it does not
import their script. The challenge is the tracker's three-target template
\nolinkurl{public/proofs/palomar-20261009/qrh/src/Challenge.lean} at the
catalogue commit of \cite{tracker}, with only the literal
$874957019421/1000000000000$ replaced by $10499/12000$ in all three
statements. The resulting \nolinkurl{evidence/kernels/src/Challenge.lean}
has SHA-256
\nolinkurl{4624fc67bbd5792408525558a48481a5cd3acd258c78caf1ec4d29528703c287}.
It was compiled with Lean 4.34.1 against an unmodified checkout of
\nolinkurl{github.com/openai/math} at \nolinkurl{fd4aeeb2e} (306 modules
of the closure of \nolinkurl{Hecke/IdealBridge} built from source) and
Mathlib \nolinkurl{d13f23b7}. The definition files
\nolinkurl{Hecke/Family.lean} and \nolinkurl{Hecke/IdealBridge.lean} of
\nolinkurl{openai/math} and \nolinkurl{DirichletContinuation.lean} and
\nolinkurl{RiemannZeta.lean} of Mathlib match their pinned git blobs
(\nolinkurl{evidence/kernels/challenge-source-audit.json}). The module
\nolinkurl{Solution} (\nolinkurl{evidence/kernels/src/Solution.lean})
imports \nolinkurl{WeightedQRH.Statements}, compiled in the Lake
workspace above, and proves the three template theorems
\nolinkurl{QRHPalomar.allDirichlet}, \nolinkurl{QRHPalomar.zeta} and
\nolinkurl{QRHPalomar.allHecke} from the endpoints. Both modules were
exported with lean4export and judged by \texttt{lake comparator} of the
official Lean v4.35.0-rc2 \nolinkurl{darwin_aarch64} release, whose
archive digest matches the digest GitHub publishes for that asset
(\nolinkurl{evidence/kernels/tool-pins.json}). After Comparator's
statement, definition and axiom checks, the external kernels con-ron
(with two jobs) and NanoDa, and then Lean's default kernel, each
replayed the export. Before judging, three controls on the same toolchain behaved as
expected (\nolinkurl{evidence/kernels/control-results.json}): a matching
pair was accepted by all three kernels; a mismatched statement was
rejected with ``Challenge and solution theorem statement do not match'';
and an ill-typed proof, a theorem whose value was replaced by its type in
the export, was rejected by con-ron, NanoDa and Lean's kernel. The judge
phase ran on 10 October 2026 from 01:44:59 to 01:53:26 UTC and exited
with status 0 (\nolinkurl{evidence/kernels/judge.log}). Its log reports
``con-ron: accepted 123717 declarations (-{}-verified)'', ``con-ron kernel
accepts the solution'', ``nanoda kernel accepts the solution'', ``Lean
default kernel accepts the solution'' and ``Your solution is okay!''. The
report \nolinkurl{evidence/kernels/result.json} has status
\nolinkurl{PASS}, verified at \nolinkurl{2026-10-10T01:53:26.235107+00:00}.
The two exports are not published; they were deleted after the run.
\nolinkurl{Challenge.export} had 382347462 bytes and SHA-256
\nolinkurl{f553d092b0690188ea806d9daadef9ec71179fd2d8e22a754dba6e40da9fbacb};
\nolinkurl{Solution.export} had 1772535340 bytes and SHA-256
\nolinkurl{7ea84555c740c3840fd7f3a3d8cdf1ec987edb60d2825425511c52e6304358a3}.
The template's zeta statement, like the zeta challenge above, has no
hypothesis $s\ne1$.

\paragraph{Earlier isolated rebuild.}
An earlier route is kept as an additional record in
\nolinkurl{evidence/isolated-rebuild/} and \nolinkurl{evidence/comparator/}.
The script \nolinkurl{verify_parallel.py} (with \nolinkurl{verify.py},
which it imports, and the configuration \nolinkurl{lean_environment.json})
rebuilt the 228 proof modules and the 3 challenge modules into a fresh
output directory, all with exit status 0, on 10 October 2026 from
00:21:47 to 00:28:17 UTC (\nolinkurl{final_rebuild_verification.json}).
The live development directory was excluded from the Lean search path,
so every local import could resolve only to an object rebuilt in that
run. The challenge modules were compiled first, against Mathlib and its
dependencies only. Mathlib, its dependencies and the original development
entered as prebuilt objects. The source hashes recorded in the Lake build
traces of those prebuilt objects of the original development were
compared with the sources at \nolinkurl{c388e88e6} and agree; that
comparison is not among the evidence files. The same object directory
also held stale objects of 253 modules deleted by the pruning; the
rebuild admitted an external import only if its module is tracked at
\nolinkurl{c388e88e6}, which excludes them. The Lake build above does not
depend on these objects. The audit log \nolinkurl{WeightedQRH.FinalAxioms.log}
of this run shows the same three axioms. With the script
\nolinkurl{run_comparator.sh} of that directory, Comparator ran against
the objects of the isolated rebuild, with its \texttt{lake build} step
replaced by a shim that only confirms that the prebuilt objects exist.
All three challenges printed ``Lean default kernel accepts the solution''
and ``Your solution is okay!''. As a negative control, the zeta
challenge with $1/2$ in place of $10499/12000$ was rejected with
``Challenge and solution theorem statement do not match''
(\nolinkurl{evidence/comparator/comparator-NegControl.log}).

\paragraph{Trust caveats.}
All these checks were run by the author on one macOS arm64 machine. They
are not a Palomar registration, an editorial review or a signed tracker
receipt, and no independent rerun has been made. The three-kernel check
used the official Lean v4.35.0-rc2 \nolinkurl{darwin_aarch64} binaries
and ran without a sandbox (\texttt{lake comparator
-{}-inadvisably-no-sandbox}); its log begins with Comparator's warning
``Sandbox disabled, this run is not trustworthy''. The tracker's
earlier kernel dossiers ran on Linux x86\_64 inside a bubblewrap
sandbox. The driver of this check was written for this run, and its
challenge module is named \nolinkurl{Challenge} rather than a random
alias. In the default-kernel Comparator runs, Comparator's Landrun
sandbox was replaced by an unsandboxed stand-in, and Comparator's
\nolinkurl{lean-toolchain} was set to v4.34.1 (pinned v4.34.0) to match
the objects; its source is otherwise unchanged. The recorded Lake build
did not exercise the dependency steps of \nolinkurl{setup.sh}, as stated
above. No external human peer review has taken place.

\paragraph{File digests.}
Every published file of the package is listed with its SHA-256 digest in
the adjacent manifest
\nolinkurl{../akashlevy-20261009-weighted-numerator.sha256.json} (tracker
path \nolinkurl{proofs/akashlevy-20261009-weighted-numerator.sha256.json}),
which lies beside the package directory rather than inside it.
Machine-specific path prefixes in the published logs and JSON reports
were replaced by neutral prefixes under \nolinkurl{/work/qrh-proof/};
\nolinkurl{publication-transformations.json} records each replacement and
the original digest of every file so changed. The digests quoted above,
of \nolinkurl{Challenge.lean} and of the two unpublished exports, do not
change on publication.

\paragraph{Reproduction.}
The certificate runs from the package root. The other commands run in
\nolinkurl{formalization/}, the last two after \nolinkurl{setup.sh}. The
package stores its files without the executable bit, so the two shell
scripts are run with \texttt{bash}. The Comparator variables name a Comparator build at the pins above, its
lean4export, and Landrun or a stand-in. For the three-kernel driver,
\nolinkurl{<pristine>} is a Lake workspace with the
\nolinkurl{lakefile.toml} and \nolinkurl{lake-manifest.json} of
\nolinkurl{evidence/kernels/pristine-workspace/}, the
\nolinkurl{lean-toolchain} of \nolinkurl{formalization/}, an unmodified
checkout of \nolinkurl{github.com/openai/math} at \nolinkurl{fd4aeeb2e} in
\nolinkurl{math/}, and the same two patches applied to its dependencies,
built by running \texttt{lake build OAI.NumberTheory.DirichletL.Hecke.IdealBridge}
in \nolinkurl{<pristine>}; \nolinkurl{<tracker>} is a checkout of the
tracker; \nolinkurl{--judge-lean} is the Lean v4.35.0-rc2 release
directory, whose \nolinkurl{bin/} holds \nolinkurl{nanoda_bin} and
\nolinkurl{con-ron}; and \nolinkurl{<fresh>} must not exist.
\par\endgroup
{\small
\begin{verbatim}
# package root: exact exponent certificate
python3 evidence/certificate/weighted_numerator_certificate.py \
    --output <certificate.json>
# formalization/: Lake build from public sources
cd formalization && bash setup.sh
# formalization/: Comparator, default Lean kernel
export COMPARATOR=<comparator> COMPARATOR_LEAN4EXPORT=<lean4export> \
    COMPARATOR_LANDRUN=<landrun>
for c in Zeta Dirichlet Hecke; do bash run_comparator.sh $c; done
# formalization/: three-kernel check of the tracker template
python3 kernels/judge_three_kernels.py --solution-workspace . \
    --challenge-workspace <pristine> \
    --template <tracker>/public/proofs/palomar-20261009/qrh/src/Challenge.lean \
    --source-lean <lean-4.34.1> --judge-lean <lean-4.35.0-rc2> \
    --lean4export <lean4export> --out <fresh>
\end{verbatim}}
{\raggedright\noindent
The certificate output coincides with
\nolinkurl{evidence/certificate/weighted_numerator_certificate.json}.
The scripts of the earlier isolated route are a record; they expect
prebuilt objects at machine-specific locations, set in
\nolinkurl{lean_environment.json} and in the scripts, which must be
adapted before reuse.\par}

\begin{thebibliography}{9}
\raggedright
\bibitem{openai} OpenAI, \emph{The Quasi-Riemann Hypothesis:
A Zero-Free Half-Plane $\Re s>7/8$}, 30 September 2026.
\nolinkurl{github.com/openai/math} commit \nolinkurl{fd4aeeb2e},
\nolinkurl{preprints/The-Quasi-Riemann-Hypothesis-September-30-2026/}.
\href{https://github.com/openai/math/blob/fd4aeeb2ee4fc729c18d98444fed42fd0529eeeb/preprints/The-Quasi-Riemann-Hypothesis-September-30-2026/paper.pdf}{Pinned manuscript}.
\bibitem{fork} OpenAI, Lean development of \cite{openai}, pruned fork by
A.~Levy: \nolinkurl{github.com/akashlevy/math}, branch
\nolinkurl{qrh-23-24-pruned}, commit
\nolinkurl{c388e88e6519d139ad2ca5571ef0c7fbf0e245b5}, derived from
\nolinkurl{github.com/openai/math} commit \nolinkurl{fd4aeeb2e}.
\href{https://github.com/akashlevy/math/tree/c388e88e6519d139ad2ca5571ef0c7fbf0e245b5}{Pinned source}.
\bibitem{mathlib} The mathlib Community, \emph{Mathlib}, commit
\nolinkurl{d13f23b723b8a846827a245b89c10fc7d3f11612}.
\href{https://github.com/leanprover-community/mathlib4/tree/d13f23b723b8a846827a245b89c10fc7d3f11612}{Pinned source}.
\bibitem{comparator} Lean FRO, \emph{Comparator}, commit
\nolinkurl{d03acab154d269c06e60e4de7e4cc85deebff94b}, with
\emph{lean4export} commit
\nolinkurl{076e8e57707e813375e8f9da8bf989799ace9680}.
\href{https://github.com/leanprover/comparator/tree/d03acab154d269c06e60e4de7e4cc85deebff94b}{Comparator};
\href{https://github.com/leanprover/lean4export/tree/076e8e57707e813375e8f9da8bf989799ace9680}{lean4export}.
\bibitem{shunia} GPT-6 Astra and J.~M. Shunia, manuscript
\nolinkurl{compensated_half_plane_refinement.pdf}, repository
\nolinkurl{jshunia/rhtm}, revision \nolinkurl{c3abed0cae8f}.
\href{https://github.com/jshunia/rhtm/blob/c3abed0cae8fdb0090c12b0e4d20aea68b6e067d/compensated_half_plane_refinement.pdf}{Pinned manuscript}.
\bibitem{guo} J.~Guo, written proof and exact-arithmetic certificates,
directory \nolinkurl{zero-free-beyond-7-8} of repository
\nolinkurl{aster2024/openai-math-followups}, revision
\nolinkurl{5750cc81}.
\href{https://github.com/aster2024/openai-math-followups/blob/5750cc81e489729a2111d6c9c073f4adecdb9de7/zero-free-beyond-7-8/README.md}{Pinned record}.
\bibitem{nielstron} Nielstron, Lean package for the algebraic endpoint,
QuasiRiemannTracker pull request~\#2, open and unmerged as of
10 October 2026, 00:16 UTC, head \nolinkurl{0284b10d}.
\href{https://github.com/GTimG/QuasiRiemannTracker/pull/2}{Pull request};
\href{https://github.com/nielstron/QuasiRiemannTracker/blob/0284b10d3c8c92b4d8995c36f24728be03815b4a/public/proofs/nielstron-algebraic-20261009-kernels/result.json}{pinned report}.
\bibitem{tracker} QuasiRiemannTracker (QRH Bounds), maintained by
T.~Gehrunger, catalogue at commit \nolinkurl{40844a76}, as of
10 October 2026, 00:16 UTC.
\href{https://gtimg.github.io/QuasiRiemannTracker/}{Website};
\href{https://github.com/GTimG/QuasiRiemannTracker/blob/40844a7614b67840801e65b836b3da2e1f7ff029/catalogue/results.json}{pinned catalogue}.
\end{thebibliography}
\end{document}
